\begin{center}
{\large\bfseries Abstract}
\end{center}

\noindent
Hands-on practice is the most effective way to learn web-application security, and
Capture-the-Flag (CTF) style exercises are the usual way to deliver it. Existing
platforms share three weaknesses in a classroom setting. They treat
every learner the same way, offering one fixed sequence of challenges regardless of
skill. They add game elements such as points and leaderboards without regard to the
well-documented ways these demotivate weaker students. And they give the
instructor little more than a solved/unsolved count, so struggling students are
noticed late. This synopsis proposes the design and development of an
\emph{adaptive, gamified web-security training platform} for penetration-testing
education. The platform hosts an intentionally vulnerable web application covering
the OWASP Top~10 vulnerability families, and scores each challenge automatically
from completion, accuracy, and time taken. A gamification module (points, badges,
levels, achievements, and a peer-banded leaderboard) is designed against the known
negative effects of gamification rather than added blindly. A hint-and-resource
layer supports beginners as they work, and a lightweight machine-learning model
uses each learner's own performance record to predict who is likely to struggle and
to recommend a suitably difficult next challenge, so that assignments adapt to the
individual. An instructor dashboard turns this record into an at-a-glance view of
progress, challenge completion, and overall performance across the cohort. The
platform is scoped as a group project, and its effectiveness will be evaluated through a pilot measuring learner
engagement and the improvement in practical cybersecurity skill.

\vspace{0.8em}
\noindent\textbf{Keywords:} web security education, penetration testing,
Capture the Flag, gamification, adaptive learning, personalization, machine
learning, learning analytics, OWASP Top~10, instructor dashboard.
